\documentclass[11pt]{article}
\usepackage{amsmath,amssymb,amsthm,enumitem,geometry}
\geometry{margin=1in}

\newtheorem{lemma}{Lemma}
\newtheorem{theorem}{Theorem}

\title{Zian Kang}
\date{}

\begin{document}

\maketitle


% ============================================================
\section*{Question 1}
\noindent Statement: Which of the following are graphic sequences? Provide a construction or a proof of impossibility for each.

\medskip


\begin{enumerate}[label=(\alph*), itemsep=10pt, leftmargin=*]
\item $(5,5,4,3,2,2,2,1)$ is graphic.

HH:
\[
(5,5,4,3,2,2,2,1)\to(4,3,2,2,1,1,1)\to(2,1,1,1,1,0)\to(1,1,0,0,0)\to(0,0,0,0).
\]
One realization:
\[
\{v_1v_2,v_1v_3,v_1v_4,v_1v_5,v_1v_6,\;
  v_2v_3,v_2v_4,v_2v_7,v_2v_5,\;
  v_3v_4,v_3v_6,\;
  v_7v_8\}.
\]

\item $(5,5,4,4,2,2,1,1)$ is graphic.

HH:
\[
(5,5,4,4,2,2,1,1)\to(4,3,3,1,1,1,1)\to(2,2,1,1,0,0)\to(1,1,0,0,0)\to(0,0,0,0).
\]
One realization:
\[
\{v_1v_2,v_1v_3,v_1v_4,v_1v_5,v_1v_6,\;
  v_2v_3,v_2v_4,v_2v_7,v_2v_8,\;
  v_3v_4,v_3v_5,\;
  v_4v_6\}.
\]

\item $(5,5,5,3,2,2,1,1)$ is graphic.

HH:
\[
(5,5,5,3,2,2,1,1)\to(4,2,1,1,1,1,0)\to(1,1,0,0,0,0)\to(0,0,0,0,0).
\]
One realization:
\[
\{v_1v_2,v_1v_3,v_1v_4,v_1v_5,v_1v_6,\;
  v_2v_3,v_2v_4,v_2v_7,v_2v_8,\;
  v_3v_4,v_3v_5,v_3v_6\}.
\]

\item $(5,5,5,4,2,1,1,1)$ is not graphic.

HH:
\[
(5,5,5,4,2,1,1,1)
\to (4,4,3,1,1,1,0)
\to (3,2,1,0,0,0)
\to (1,0,-1,0,0).
\]
In the last step, deleting $3$ from $(3,2,1,0,0,0)$ and subtracting $1$ from the next three terms gives
$(2,1,0,0,0)\mapsto(1,0,-1,0,0)$, which has a negative entry. Hence the sequence is not graphic.
\end{enumerate}

\newpage
% ============================================================
\section*{Question 2}
\noindent Statement: Let $d=(d_1,\dots,d_n)$ be nonincreasing. Fix $k$ and form $d'$ by deleting $d_k$
and subtracting $1$ from the $d_k$ largest remaining entries. Prove that $d$ is graphic iff $d'$ is graphic.

\begin{proof}
~
\begin{enumerate}
    \item Write the entries of $d$ as degrees assigned to labeled vertices $v_1,\dots,v_n$,
with $\deg(v_i)=d_i$. Let $x=v_k$ and $r=d_k$, fixing a $k$.
Let $A$ be a set of $r$ vertices among $V\setminus\{x\}$ whose degrees are the $r$ largest among
$\{d_i:i\neq k\}$.

\item Switching lemma:
For every $x,y,z,w$ distinct, if $xy$ and $zw$ are edges, and $xz$ and $yw$ are nonedges,
then replacing $xy,zw$ by $xz,yw$ preserves all degrees and keeps the graph simple.

\item $\Rightarrow$ Assume $d$ is graphic; let $G$ be a simple graph with degree sequence $d$.
Claim there is a realization $G^*$ of $d$ in which $x$ is adjacent to every vertex in $A$.

If $x$ already hits all of $A$, done. Otherwise pick $z\in A\setminus N_G(x)$.
Since $z\notin N_G(x)$, there exists some neighbor $y\in N_G(x)\setminus A$.
By choice of $A$, we have $\deg(z)\ge \deg(y)$.

If $N_G(z)\subseteq N_G(y)$, then $\deg(z)\le \deg(y)$, forcing $\deg(z)=\deg(y)$; but then we may swap
membership of $y$ and $z$ in $A$, contradicting $z\in A\setminus N_G(x)$ and $y\notin A$.
So $N_G(z)\not\subseteq N_G(y)$, choose $w\in N_G(z)\setminus N_G(y)$.
Then $zw$ is an edge, $yw$ is a nonedge, and also $xz$ is a nonedge by construction, while $xy$ is an edge.
Apply the switching lemma to edges $xy,zw$ to get a new realization where $x$ is adjacent to $z$.
This increases $|A\cap N(x)|$. Repeating yields a realization $G^*$ where
$A\subseteq N_{G^*}(x)$.

Now delete $x$ from $G^*$. The degrees of the vertices in $A$ drop by $1$, and all other degrees stay the same.
This is exactly the sequence $d'$. Thus $d'$ is graphic.

\item $\Leftarrow$ Assume $d'$ is graphic, realized by a simple graph $H$ on $n-1$ vertices.
Take the $r$ vertices of $H$ corresponding to the $r$ largest entries (before sorting) in $d'$, and add a new vertex
$x$ adjacent to those $r$ vertices. This increases their degrees by $1$ and gives $\deg(x)=r$,
producing degree sequence $d$. Hence $d$ is graphic.
\end{enumerate}
\end{proof}

\newpage
% ============================================================
\section*{Question 3}
\noindent Statement: Let $d_1\ge\cdots\ge d_n\ge 0$. Prove there exists a loopless graph allowing multiple edges
(a loopless multigraph) with degree sequence $d_1,\dots,d_n$ iff $\sum_i d_i$ is even and $d_1\le d_2+\cdots+d_n$.

\begin{proof}
~
\begin{enumerate}
    \item $\Rightarrow$ In any (multi)graph, $\sum_i d_i=2|E|$ is even by handshaking.
Also, each edge incident to the vertex of degree $d_1$ uses $1$ unit of degree from some other vertex.
Thus the total available degree outside that vertex is $d_2+\cdots+d_n$, so we must have $d_1\le d_2+\cdots+d_n$.

\item $\Leftarrow$ Assume $\sum_i d_i$ is even and $d_1\le d_2+\cdots+d_n$.
Build a loopless multigraph by repeatedly adding edges between distinct vertices with positive remaining degree.

Maintain a list of remaining degrees. While some degree is positive, pick a vertex $u$ of maximum remaining degree $m$.
Since $m\le$ (sum of the other remaining degrees), there exists another vertex $v\neq u$ with positive remaining degree.
Add an edge $uv$ (with possibility to add multiple edges), and decrease the remaining degrees of
$u$ and $v$ by $1$.

This never creates a loop since endpoints are distinct, never makes degrees negative, and decreases the total remaining sum by $2$,
so it must terminate. If it were to get stuck with exactly one vertex having positive remaining degree, then we would have
$M>0$ but sum of others $=0$, contradicting the maintained inequality $M\le$ (sum of others).
Hence the process ends with all remaining degrees $0$, producing the desired loopless multigraph.
\end{enumerate}
\end{proof}

\newpage
% ============================================================
\section*{Question 4}
\noindent Statement: Let $d=(d_1,\dots,d_{2k})$ be defined by $d_{2i-1}=d_{2i}=i$ for $1\le i\le k$. Prove $d$ is graphic.

\begin{proof}
Construct a simple bipartite graph with bipartition
\[
A=\{a_1,\dots,a_k\},\qquad B=\{b_1,\dots,b_k\},
\]
and edges
\[
a_i b_j \in E \quad\text{iff}\quad j\le i.
\]
Then $\deg(a_i)=i$ (neighbors $b_1,\dots,b_i$).
Also $\deg(b_j)=|\{i: j\le i\le k\}|=k-j+1$.

Thus the multiset of degrees is
\[
\{\deg(a_1),\dots,\deg(a_k)\}\cup\{\deg(b_1),\dots,\deg(b_k)\}
=\{1,2,\dots,k\}\cup\{1,2,\dots,k\},
\]
which is exactly $d$ (up to reordering). Hence $d$ is graphic.
\end{proof}

\newpage
% ============================================================
\section*{Question 5}
\noindent Statement: Prove that a digraph is strongly connected iff for each partition $V=S\cup T$ into nonempty sets,
there is an edge from $S$ to $T$.

\begin{proof}
~
\begin{enumerate}
    \item $\Rightarrow$ If $D$ is strongly connected, pick any nonempty partition $V=S\cup T$ and choose $s\in S$, $t\in T$.
There is a directed $s\to t$ path. Along that path, let $uv$ be the first edge with $u\in S$ and $v\in T$.
Then $uv$ is an edge from $S$ to $T$.

\item $\Leftarrow$ Assume: for every partition $V=S\cup T$ into nonempty sets,
there exists an edge from $S$ to $T$.

Fix any vertex $s\in V$. Let
\[
R=\{v\in V:\text{ there is a directed path from } s \text{ to } v\}
\]
be the set of vertices reachable from $s$. Clearly $s\in R$, so $R\neq\emptyset$.

If $R\neq V$, then $T:=V\setminus R$ is nonempty, and there is a partition $V=R\cup T$ into nonempty sets.
By the assumption, there is an edge $u\to w$ with $u\in R$ and $w\in T$.
But then $w$ is reachable from $s$, so $w\in R$,
contradicting $w\in T=V\setminus R$. Hence $R=V$.

Therefore, from this arbitrary $s$, every vertex is reachable. Since $s$ is arbitrary, every vertex reaches every other,
so $D$ is strongly connected.
\end{enumerate}
\end{proof}

\newpage
% ============================================================
\section*{Question 6}
\noindent Statement: For each $n\ge 7$, prove or disprove: every simple digraph on $n$ vertices has two vertices
with the same out-degree or two vertices with the same in-degree.

\begin{proof}
\emph{Disprove.} Construct, for every $n\ge 2$ (hence for every $n\ge 7$), a simple digraph in which
all out-degrees are distinct and all in-degrees are distinct.

Let $V=\{v_1,\dots,v_n\}$. For each pair $i<j$, put exactly one directed edge, namely
\[
v_i \to v_j.
\]
This is a tournament, hence a simple digraph.

Now compute degrees. For a fixed $i$:
\begin{itemize}
\item $v_i$ has an out-edge to each $v_j$ with $j>i$, so $d^+(v_i)=|\{j:i<j\le n\}|=n-i$.
\item $v_i$ has an in-edge from each $v_j$ with $j<i$, so $d^-(v_i)=|\{j:1\le j<i\}|=i-1$.
\end{itemize}

Thus
\[
\{d^+(v_1),\dots,d^+(v_n)\}=\{n-1,n-2,\dots,0\},
\qquad
\{d^-(v_1),\dots,d^-(v_n)\}=\{0,1,\dots,n-1\},
\]
so all out-degrees are pairwise distinct and all in-degrees are pairwise distinct. Therefore the claimed statement is false.
\end{proof}

\newpage
% ============================================================
\section*{Question 7}
\noindent Statement: Let $G$ be an $n$-vertex digraph with no directed cycles. Prove its vertices can be ordered
$v_1,\dots,v_n$ so that if $v_iv_j\in E(G)$ then $i<j$.

\begin{proof}
Prove the claim by induction on $n=|V(G)|$.

\medskip
\noindent Claim: $G$ has a vertex of indegree $0$.

\begin{proof}
    Suppose not; then every vertex has indegree at least $1$.
Pick any vertex $x_1$. Since $d^-(x_1)\ge 1$, choose an in-neighbor $x_2$ with $x_2\to x_1$.
Similarly choose $x_3$ with $x_3\to x_2$, and continue. This produces a sequence
$x_1,x_2,x_3,\dots$ of vertices in a finite set, so some vertex repeats: $x_i=x_j$ with $i<j$.
Then
\[
x_j \to x_{j-1} \to \cdots \to x_{i+1} \to x_i=x_j
\]
is a directed cycle, contradicting that $G$ is acyclic. Therefore some vertex has indegree $0$.
\end{proof}

\medskip
\noindent Proof by Induction:

If $n=1$, any ordering works.

Assume $n\ge 2$ and the statement holds for all smaller acyclic digraphs.
By Claim 1, pick a vertex $v_1$ with indegree $0$. Let $G' = G - v_1$ be the digraph geted by deleting $v_1$
and all incident edges. Then $G'$ is still acyclic.

By the induction hypothesis, $G'$ has an ordering $v_2,\dots,v_n$ such that every edge in $G'$
goes from earlier to later in this list. Then place $v_1$ to get the order $v_1,v_2,\dots,v_n$.

Now check edges in $G$:
\begin{itemize}
\item Any edge not incident to $v_1$ lies in $G'$, so it goes forward.
\item There is no edge into $v_1$ (since $d^-(v_1)=0$), so we never have an edge going from a later vertex to $v_1$.
\item Any edge out of $v_1$ goes from $v_1$ to some $v_j$ with $j\ge 2$, hence also goes forward.
\end{itemize}
Therefore, in the order $v_1,\dots,v_n$, every directed edge goes from a smaller index to a larger index.
\end{proof}

\newpage
% ============================================================
\section*{Question 8}
\noindent Statement: Use Lemma 1.4.23 and induction on the number of edges to prove the characterization of Eulerian digraphs (Theorem 1.4.24).
\begin{proof}
~
\begin{enumerate}
    \item $\Rightarrow$ Suppose $D$ has an Eulerian directed circuit $W$.
Each passage of $W$ through a vertex uses one incoming edge and one outgoing edge. Since $W$ is closed,
the first edge used at a vertex is paired with the last edge used there as well, so the total number of times
$W$ enters $v$ equals the total number of times it leaves $v$. Hence $d^+(v)=d^-(v)$ for every $v$.

Also, two edges can lie in the same directed trail only when they lie in the same weak component of the digraph. Since $W$ contains every edge, all edges lie in one
weak component. Thus $D$ has at most one nontrivial weak component.

\item $\Leftarrow$ Construct by induction on the number of edges $m=|E(D)|$.

\begin{enumerate}
    \item Basis step: $m=0$.
A closed directed trail consisting of a single vertex is an Eulerian circuit.

\item Induction step: $m>0$.
Let $D$ have $m$ edges. Let $H$ be the nontrivial weak component of $D$.
For every vertex $v$ of $H$, we have $d^+(v)+d^-(v)\ge 1$, and this implies $d^+(v)\ge 1$. Hence $\delta^+(H)\ge 1$.

By Lemma 1.4.23, $H$ contains a directed cycle $C$.
Let $D' = D - E(C)$ be the digraph geted by deleting the edges of $C$.

Since $C$ uses either $0$ or $1$ incoming edge and the same number of outgoing edges at each vertex,
deleting $E(C)$ decreases $d^+$ and $d^-$ equally at every vertex. Therefore every vertex of $D'$ still satisfies
$d^+(v)=d^-(v)$.

Now each weak component of $D'$ is balanced and has fewer than $m$ edges, so use the induction hypothesis
to conclude that each nontrivial weak component of $D'$ has an Eulerian directed circuit.

To combine these into an Eulerian circuit of $D$, traverse the directed cycle $C$. When a nontrivial component of
$D'$ is entered for the first time, let it go along an Eulerian circuit of that
component and return to the same vertex as the beginning. Continuing this process for every nontrivial component of
$D'$, we get a closed directed trail that uses every edge of $C$ exactly once
and every edge of $D'$ exactly once. Hence we have constructed an Eulerian directed circuit of $D$.
\end{enumerate}
\end{enumerate}
\end{proof}

\newpage
% ============================================================
\section*{Question 9}
\noindent Statement: By Proposition 1.4.30, every tournament has a king. Let $T$ be a tournament having no vertex with indegree $0$.
\begin{enumerate}[label=(\alph*)]
\item Prove that if $x$ is a king in $T$, then $T$ has another king in $N^-(x)$.
\item Use part (a) to prove that $T$ has at least three kings.
\item For each $n\ge 3$, construct a tournament $T$ with $\delta^-(T)>0$ and only $3$ kings.
\end{enumerate}

\medskip
\begin{proof}
~
\begin{enumerate}[label=\alph*]
    \item Since $\delta^-(T)>0$, $N^-(x)\neq\emptyset$. Consider the subtournament $T[N^-(x)]$.
By the fact that every tournament has a king, $T[N^-(x)]$ has a king $z$.

Let $v\in V(T)$ be arbitrary.
If $v\in N^-(x)$, then $z$ reaches $v$ in $\le 2$ steps inside $T[N^-(x)]$.
If $v=x$, then $z\to x$ since $z\in N^-(x)$.
If $v\in N^+(x)$, then $z\to x\to v$.
Thus $z$ reaches every $v$ within $2$ steps, so $z$ is a king in $T$ and lies in $N^-(x)$.
Also $z\neq x$, giving another king.
\item Pick any king $x$. By (a) there is a king $y\in N^-(x)$, hence $y\neq x$.
Apply (a) to $y$ to get a king $z\in N^-(y)$. Then $z\neq y$.
Also $z\neq x$ because $y\in N^-(x)$ means $y\to x$, hence $x\notin N^-(y)$.
So $x,y,z$ are three distinct kings.

\item Construction for each $n\ge 3$.
Let vertices be $\{a,b,c\}\cup\{u_1,\dots,u_{n-3}\}$.
Orient $a\to b$, $b\to c$, $c\to a$, and for each $i$ orient
\[
a\to u_i,\quad b\to u_i,\quad c\to u_i.
\]
Orient edges among the $u_i$ arbitrarily.

Then $\delta^-(T)>0$ because $a,b,c$ each have indegree $1$ from the $3$-cycle and each $u_i$ has indegree $\ge 3$.
Vertices $a,b,c$ are kings: each reaches every $u_i$ in one step, and reaches the third vertex among $a,b,c$ in two steps.
No $u_i$ is a king since $u_i$ has no directed path to $a$ because all edges between
$\{a,b,c\}$ and $\{u_i\}$ point from $\{a,b,c\}$ to $\{u_i\}$.
Hence exactly three kings.
\end{enumerate}
\end{proof}

\newpage
% ============================================================
\section*{Question 10}
\noindent Statement: Consider the following algorithm whose input is a tournament $T$.
\begin{enumerate}[label=(\arabic*)]
\item Select a vertex $x$ in $T$.
\item If $x$ has indegree $0$, call $x$ a king of $T$ and stop.
\item Otherwise, delete $\{x\}\cup N^+(x)$ from $T$ to form $T'$.
\item Run the algorithm on $T'$; call the output a king in $T$ and stop.
\end{enumerate}
Prove that this algorithm terminates and produces a king in $T$.

\medskip
\begin{proof}
~
\begin{enumerate}
    \item Termination
    
    Each recursive call deletes at least the chosen vertex $x$,
so the number of vertices strictly decreases. Eventually one vertex remains; it has indegree $0$,
so the algorithm stops.

\item Correctness

Prove by induction on $n=|V(T)|$ that the algorithm outputs a king.

Base case $n=1$ is trivial. For $n>1$, pick $x$.
If $d^-(x)=0$, then $x$ beats every other vertex, so reaches everyone in one step; it is a king.

Otherwise $N^-(x)\neq\emptyset$. The remaining vertex set after deleting $\{x\}\cup N^+(x)$ is exactly $N^-(x)$,
so the recursive call runs on the subtournament $T'=T[N^-(x)]$.
By induction the output $y$ is a king of $T'$.

Let $v\in V(T)$.
If $v\in N^-(x)$, then $v\in V(T')$ and $y$ reaches $v$ in $\le 2$ steps in $T'$.
If $v=x$, then $y\to x$ since $y\in N^-(x)$.
If $v\in N^+(x)$, then $y\to x\to v$.
Thus every $v$ is within directed distance $\le 2$ from $y$, so $y$ is a king of $T$.
\end{enumerate}
\end{proof}

\end{document}